\originalsection{Airy转折点与两端匹配}
右转折点$x_2$若$V'(x_2)=F>0$, 局部$V-E\simeq F(x-x_2)$. 定义长度$d=[\hbar^2/(2mF)]^{1/3}$及$z=(x-x_2)/d$, 方程成为$u_{zz}=zu$, 称Airy方程. $z>0$是右侧禁止区,$z<0$是左侧允许区. 实解$\operatorname{Ai}(z)$在正轴衰减,$\operatorname{Bi}(z)$增长; 以下推导这两个主阶与振荡端的关系, 不将它们当成未经说明的相位补丁.

一个可验证的积分定义是$\operatorname{Ai}(z)=(2\pi\ii)^{-1}\int_{\mathcal C}\exp(s^3/3-zs)\dd s$, 路径从辐角$-\pi/3$的无穷走到$+\pi/3$的无穷, 两端指数实部负. 对$z$求两次导数, 积分$\partial_s\exp(s^3/3-zs)=(s^2-z)\exp(s^3/3-zs)$两端消失, 即$\operatorname{Ai}''-z\operatorname{Ai}=0$. 这一积分工具涉及复路径积分; 下述变形使用整函数可沿衰减路径变形的Cauchy定理, 作为进阶数学输入, 无需为量子入门先修整套复分析.

当$z\gg1$, 衰减鞍点$s=\sqrt z$, 取其附近$s=\sqrt z+\ii y$, 指数恰为$-2z^{3/2}/3-\sqrt z\,y^2-\ii y^3/3$. 有效$y$范围$z^{-1/4}$, 三次项是小量; Gaussian积分给
\[
\operatorname{Ai}(z)\sim\frac1{2\sqrt\pi}z^{-1/4}\e^{-2z^{3/2}/3}.
\]
路径变形两端仍留在衰减扇区, 无穷端没有额外边界贡献. 不是把一条会指数爆炸的任意半圆路径强行闭合.

负轴令$z=-q,q\gg1$, 同一积分可变到振荡表示$\operatorname{Ai}(-q)=\pi^{-1}\int_0^\infty\cos(t^3/3-qt)\dd t$, 以振荡积分极限理解. 相位在$t_0=\sqrt q$驻定, 其值$-2q^{3/2}/3$, 二阶项$\sqrt q(t-t_0)^2$. 以Gaussian正则化定义Fresnel积分$\int_{-\infty}^\infty\e^{\ii a y^2}\dd y=\e^{\ii\pi/4}\sqrt{\pi/a}$,$a>0$, 得
\[
\operatorname{Ai}(-q)\sim\frac1{\sqrt\pi}q^{-1/4}\sin(2q^{3/2}/3+\pi/4).
\]
驻定区扩到全线只取主阶, 非驻定区可分部积分控制; 准确地说$\operatorname{Ai}(-q)=\pi^{-1/2}q^{-1/4}[\sin(2q^{3/2}/3+\pi/4)+O(q^{-3/2})]$, 误差按包络计; 主振荡项的零点附近不能报告统一相对误差. $\pi/4$来自振荡Gaussian的相位, 系数2来自衰减端与振荡端的$1/(2\sqrt\pi)$、$1/\sqrt\pi$差, 都有具体计算来源.

另一实解可定义为$\operatorname{Bi}(z)=\pi^{-1}\int_0^\infty[\e^{-t^3/3+zt}+\sin(t^3/3+zt)]\dd t$. 对$z$二次求导, 两部分分部积分的端点项互相抵消, 同样满足Airy方程. 正$z$时第一积分在$t=\sqrt z$有实最大点, Laplace的Gaussian计算给$\operatorname{Bi}(z)\sim\pi^{-1/2}z^{-1/4}\e^{2z^{3/2}/3}$. 负$z=-q$时, 非驻定端点贡献在两项间抵消; 正驻定点的正弦Gaussian给$\operatorname{Bi}(-q)\sim\pi^{-1/2}q^{-1/4}\cos(2q^{3/2}/3+\pi/4)$. 这两组渐近组成独立的实解基, 不需假定禁止区只有衰减一支才能匹配一般解.

在线性势的重叠区,$|p|=\hbar\sqrt{|z|}/d$, $\int|p|\dd x/\hbar=2|z|^{3/2}/3$. 与上述幅因子对照, 得两个连接:
\[
\begin{aligned}
\frac C{\sqrt{|p|}}\e^{-\int_{x_2}^x|p|\dd s/\hbar}
&\longleftrightarrow\frac{2C}{\sqrt p}\sin\left(\int_x^{x_2}p\dd s/\hbar+\frac\pi4\right),\\
\frac D{\sqrt{|p|}}\e^{+\int_{x_2}^x|p|\dd s/\hbar}
&\longleftrightarrow\frac D{\sqrt p}\cos\left(\int_x^{x_2}p\dd s/\hbar+\frac\pi4\right).
\end{aligned}
\]
要求有一段既能线性化势、又有$|z|\gg1$允许WKB的重叠区. 对左简单转折点把坐标方向反转, 得相应连接. 硬墙、高阶转折点或两个几乎合并的转折点须不同匹配.

对两端皆衰减束缚态, 令$A(x)=\int_{x_1}^xp\dd s/\hbar$, $I=\int_{x_1}^{x_2}p\dd s/\hbar$. 左匹配要求允许解正比$\sin(A+\pi/4)$, 右匹配要求正比$\sin(I-A+\pi/4)$. 展成$\sin A,\cos A$比较两系数, 要两者成比例必须$\sin(I+\pi/2)=0$, 即$I=\pi(n+1/2)$,$n=0,1,\ldots$. 这是两端局部相位共同筛能量, 并非在经典作用量上任意加半个量子数.
